% !TeX root = ../thesis.tex


%The first chapter of your thesis.
\chapter{introduction}\label{ch:introduction}



% Placeholder prose for previewing the campus layout; replace with your own text.
\section{Asian hornet}

The Asian hornet forms a problem for both wildlife and human health. Therefore, a cost-effective and easy-to-use tracker is needed to find and remove nests that are a risk to apiaries or human health. Hornet tracking can be split up into two main categories: visual tracking and radio tracking. In visual tracking, the nest is found by visually following the hornet or writing down the flight direction, while radio tracking relies on catching a hornet and applying a tag that can be passive and work with reflection, or that is active and sends out its own radio waves. This thesis focuses on making a radio telemetry tag; this is the active kind. 

\section{Thesis objective}

The objective of this thesis is to design a cost-effective, rechargeable, sub-$200\,\text{mg}$ transmitter that is practical and user-friendly for non-technical operators. The main challenge will be to make a miniature design without losing this usability. This will be achieved by following the primary objectives.

\section{Primary objectives}

\subsection{Literature review}
Here the objective is to learn about the current and past research in the field of radio telemetry and Asian hornets. The commercial trackers will also be compared to look at what the industry is using at the moment. With the gathered information, a design can be chosen that fits this thesis the best.

\subsection{Theoretical model and simulations} 

With the design decision made, a theoretical model now needs to be built. This serves the function of understanding how the transmitter works and how the components need to be calculated. The simulation will show if the design and component calculations are correct.

\subsection{Develop initial prototype}

This prototype will validate the simulations and can be experimented with. This prototype will not be miniature. 

\subsection{Optimize the transmitter for hornets}

The objective of this part will be to scale down the transmitter to sub-$200\,\text{mg}$ and make the tag attachable to hornets.

\subsection{Lab testing}

This part aims to validate the range, frequency stability, battery life, and fine-tune where needed. 

\subsection{Field testing}

At this point, the tag is ready to be tested in the field. The main objective is to validate the hornet attachment mechanism and to see how well hornets cope with flying with the tag.

%\subsection{Conclusions}